\subsection{The completed sixteen-arrow linear calculation}

\begin{theorem}\label{C:sixteen-three-linear}
No three-object finite category with linear acyclic support
and at most sixteen arrows has classifying space homotopy
equivalent to $\mathbb CP^2$.
\end{theorem}
\begin{proof}
Theorem~\ref{C:CP2-lower16} handles smaller arrow totals.
At total sixteen, write the corner orders as $(m,n,p)$
and the forward hom-set sizes as $(x,y,z)$ in directions
$01,12,02$.  All six integers are positive, and their sum
is sixteen.  There are $\binom{15}{5}=3003$ such ordered
six-tuples.  Opposition and object relabeling give exactly
$1547$ canonical Hom-cardinality profiles.

The structural omissions use the singleton joins and
singleton cofinal deletions, the fixed-arrow lemma, the
two-block zero-product lemma, and the triangle-band
compression already proved above.  Large middle corners
and the necessary extreme-corner catalog are treated in
Proposition~\ref{C:sixteen-linear-large-corners}.
The later exact middle-algebra reductions are
Corollary~\ref{C:semisimple-ends-middle-four} and
Corollary~\ref{C:semisimple-ends-binary-homs}.
The final endpoint reduction is
Theorem~\ref{C:semisimple-endpoints-sixteen}, with its
stated initial-left and final-right projectivity hypotheses.
No corresponding assertion for arbitrary larger arrow
totals is used.

The search was continued through successive completed-profile
prefixes as these reductions became available.  Across all
six retained streams, exactly $347$ distinct canonical
profiles have complete table searches.  The other $1200$
canonical profiles have one of the proved structural
obstructions.  A separate inventory enumerates the six
positive sizes directly, rather than the producer's
nine-bin distribution of arrows.  It reconstructs every
canonical profile and every structural reason, checks
that all completed profile sets are disjoint, and matches
the final required order of the remaining profiles.
Nontransitive assignments with two nonempty adjacent
hom-sets and an empty composite hom-set cannot be categories
and are not counted as valid linear profiles.

For each retained profile, all corner monoids of order at
most seven range over the complete ordinary catalogs,
including both orientations.  Only the proved possible-
positive-$H^2$ catalog is needed for an extreme corner of
order eight.  It contains $1032$ distinct literal table
orientations.  The number of literal corner tables across
the eight order catalogs is $36459$.
Each remaining cross product ranges over its entire typed
target hom-set.  Associativity propagation imposes only
equalities forced by the category axioms, and every
unforced value is tried.  Thus the search retains every
category in its stated scope.  Exact permutations inside
cross hom-sets may remove repeated tables; no relative
position of the corner actions is omitted.

The endpoint flags used in the final stream are separately
certified.  For every one of the $36459$ literal tables,
\path{code/b_verify_linear16_projective_endpoints.py}
solves the actual rational invariant-vector equations
on both sides, with augmentation equal to one.  It
compares these exact criteria with the producer's
minimal-ideal cardinality flags.  All $72918$ comparisons
agree, and all $12320569$ corner associativity equations
are checked.  The result is
\path{code/outputs/b_linear16_projective_endpoints_verified.json}.
The independent profile inventory also predicts the two
new producer counters exactly: $5311$ early endpoint
profile omissions and $838045$ projective-endpoint corner
choices.  Both equal the actual final footer.

For every completed table, an initial normalized-bar filter
retains the necessary mod-two conditions $H_1=0$ and
$\dim H_2=1$.  The independent normalized-bar computation
\path{code/c_category_tsv16_cups.cpp} then reconstructs
the literal category and all normalized chains through
degree four.  It checks the typed category axioms and
identifies a nonboundary degree-two cocycle.  Either it
supplies and evaluates a degree-three primitive for the
Alexander--Whitney square, or it independently computes
positive third mod-two homology.  The mutually exclusive
completed results are
\[
\begin{array}{c|r|r|r|r}
 \text{stream}&\text{complete profiles}&\text{retained rows}
   &u^2=0&H_3(\mathbb F_2)\ne0\\\hline
 \text{original}&65&1401244&1347792&53452\\
 v2&51&723515&709304&14211\\
 v3&59&865106&855934&9172\\
 v4&107&1592464&1585896&6568\\
 v5&10&415822&406100&9722\\
 v6&55&1080989&1069881&11108\\\hline
 \text{total}&347&6079140&5974907&104233.
\end{array}
\]
There is no surviving candidate.  Every retained row
therefore has an obstruction to the homology or cup ring
of $\mathbb CP^2$, proving the theorem.
\end{proof}

\begin{remark}[Completed prefixes and full scope]
\label{C:sixteen-linear-complete-verification}
The final producer is
\begin{quote}\small
\path{code/c_categories16_linear_resume_v6.cpp}.
\end{quote}
Its actual final footer records $97067056$ finished
corner-representative tables and $1080989$ retained rows
in that stream.  Its reported $292$ previously completed
keys agree with the disjoint prior sets of sizes
$65,51,59,107,10$; all $55$ newly required profiles finish.
The producer's larger distribution counter also includes
invalid supports rejected before a profile starts and
is not substituted for the completed-profile count.

Interrupted historical streams remain unchanged.
Only their exact completed-profile prefixes in the table
above are used.  The original source's storage-failure
tail and every later unfinished profile are excluded.
The bounded orbit cache uses an injective packed key
for the typed multiplication entries; clearing a full
cache can retain additional repetitions and cannot
discard a previously unseen category.

The separate completed-prefix gates match every literal
source, target and multiplication table with its saved
cochain record, require the verifier's final count and
the compressed certificate footer, and hash each exact
source prefix.  The final program
\path{code/b_finish_sixteen_linear_scope.py}
then re-reads all six prefixes.  It independently checks
every row's Hom profile, every profile boundary, all
six source hashes, all certificate-interval unions,
and the actual final source EOF.  It requires the full
inventory gate
\path{code/outputs/b_sixteen_linear_inventory_complete.json};
a stopped or merely growing stream cannot satisfy it.

The complete checks comprise $2767269079$ typed
associativity equations, $1483020541$ cocycle equations,
and $6463509055$ literal square-primitive equations.
Their final record is
\begin{quote}\small
\path{code/outputs/b_sixteen_linear_complete_verified.json},
\end{quote}
with status
\begin{quote}\small\ttfamily
PASS\_COMPLETE\_SIXTEEN\_ARROW\_LINEAR\_SCOPE.
\end{quote}
This records the full independent table-to-bar computation
and exact source linkage; no uncomputed homology tail is
used in the exclusion.
\end{remark}

\begin{theorem}[The global lower bound of seventeen arrows]
\label{C:CP2-lower17}
Every finite category whose classifying space is homotopy
equivalent to $\mathbb CP^2$ has at least seventeen arrows,
identities counted.  Together with the explicit construction,
the minimum-arrow interval is $17$--$22$.
\end{theorem}
\begin{proof}
Theorem~\ref{C:CP2-lower16} excludes all smaller totals,
so it remains to exclude sixteen arrows.  Connectedness
is necessary.  The one- and two-object cases are covered
by Corollary~\ref{C:sixteen-two-objects}.
For three objects, a nontrivial strongly connected
component is excluded by
Theorem~\ref{C:sixteen-three-SCC}.  An acyclic connected
support on three objects is an inward or outward star,
or a transitive linear support: a composable pair of
strict arrows forces its long composite hom-set to be
nonempty.  Apply Theorem~\ref{C:sixteen-all-stars} and
Theorem~\ref{C:sixteen-three-linear}, respectively.
Four objects are excluded by
Theorem~\ref{C:sixteen-four-objects}, five and six by
Theorem~\ref{C:sixteen-five-six-objects}, and every larger
object count by Theorem~\ref{C:sixteen-seven-objects}.
These cases exhaust all finite categories with sixteen
arrows.  The upper bound is Theorem~\ref{C:CP2-upper22}.
\end{proof}
